% =============================================================================
%    CHAPTER 2: BACKGROUND STUDY AND LITERATURE REVIEW
% =============================================================================
\chapter{Background Study and Literature Review}

\section{Background Study}

Traditional self-supervised learning approaches in computer vision have relied either on generative pixel reconstruction or on contrastive learning across augmented views. Pixel reconstruction forces models to spend representational capacity on high-frequency, unessential visual details, while contrastive learning requires large batch sizes and carefully designed negative sampling. To address these limitations, Joint Embedding Predictive Architectures (JEPA) perform non-generative prediction directly within a learned latent feature space. \citet{assran2023ijepa} introduced Image-JEPA (I-JEPA), demonstrating that predicting missing target representations from context patches in latent space produces highly semantically rich visual embeddings.

Extending the JEPA paradigm to temporal data, \citet{bardes2024vjepa} proposed V-JEPA, which applies feature prediction to spatiotemporal video blocks. Subsequently, \citet{assran2025vjepa2} developed \vjepaa{}, scaling video latent prediction using a ViT-Huge backbone (654M parameters) trained on extensive unlabelled video collections. Most recently, \citet{murlabadia2026vjepa21} introduced V-JEPA 2.1, incorporating a dense predictive loss and multi-layer intermediate supervision to enhance spatial and temporal feature grounding. Variational extensions to JEPA have also been explored by \citet{huang2026vjepa}, who proposed Variational JEPA (VJEPA), introducing a probabilistic latent prediction framework that models distribution uncertainty for stochastic control and predictive world modeling.

Parallel to self-supervised feature learning, Vision-Language Models (VLMs) have emerged as powerful general-purpose foundation models. As surveyed by \citet{bordes2024vlm} and \citet{zhang2023vlmsurvey}, VLMs connect visual encoders with autoregressive language model decoders through visual instruction tuning. Recent developments in vision-language architectures include the \qwenvl{} technical report \citep{bai2025qwen2vl} and Qwen3-VL \citep{qwen2025qwen3vl}, which support native spatiotemporal visual tokenization, variable frame resolution, and long-context multimodal understanding. To enable practical deployment under strict hardware resource boundaries, \citet{pathak2026efficientvlms} reviewed quantization, pruning, and low-bit inference techniques for large vision-language models.

\section{Literature Review}


The progression from Image-JEPA to V-JEPA~2.1 reveals a systematic trajectory of architectural refinement aimed at producing increasingly expressive latent representations. \citet{assran2023ijepa} first established the foundational principle that predicting target patch representations from context patches---without pixel-level reconstruction---produces semantically richer embeddings than generative or contrastive alternatives. Their Image-JEPA demonstrated that a non-generative latent predictive objective avoids wasting model capacity on perceptually irrelevant high-frequency detail, while simultaneously eliminating the dependency on hand-crafted data augmentations inherent to contrastive learning pipelines.

Building on this image-level foundation, \citet{bardes2024vjepa} extended the joint-embedding predictive paradigm to the spatiotemporal domain with V-JEPA, introducing masked spatiotemporal block prediction in video feature space. The critical insight was that temporal coherence across video frames provides an especially rich supervisory signal for latent prediction, as adjacent frames share semantic content while exhibiting continuous motion dynamics. \citet{assran2025vjepa2} then scaled the V-JEPA framework significantly with V-JEPA~2, employing a ViT-Huge backbone with 654 million parameters trained on large-scale unlabelled video data. This scaling demonstrated that increasing both model capacity and training data volume within the JEPA paradigm yields representations that transfer competitively across diverse downstream tasks.

The most recent architectural advancement, V-JEPA~2.1 \citep{murlabadia2026vjepa21}, introduced two critical innovations: a dense per-token prediction loss that supervises each masked spatiotemporal token individually, and multi-layer intermediate supervision applied across encoder depth. These modifications fundamentally change the nature of learned representations by ensuring that intermediate encoder layers also acquire strong semantic features, rather than concentrating representational quality only in the final layer. This property is especially relevant to the present study, as it motivates the extraction of multi-layer features for the proposed HAF-STD Net architecture. In parallel, \citet{huang2026vjepa} proposed Variational JEPA, which reformulates latent prediction as a probabilistic inference problem by introducing distribution uncertainty modelling, demonstrating that JEPA-style architectures can be extended beyond deterministic prediction to support stochastic control and predictive world modeling.


Integrating JEPA principles into multimodal and language domains represents an active frontier of research. \citet{chen2026vljepa} proposed VL-JEPA, extending joint-embedding predictive learning to multimodal vision-language alignment without relying on autoregressive token generation. By performing cross-modal latent prediction between visual and textual representations, VL-JEPA demonstrated that the non-generative JEPA objective can bridge modality gaps while preserving the computational advantages of avoiding pixel or token reconstruction. \citet{huang2025llmjepa} combined Large Language Models with JEPA, demonstrating latent-space predictive loss objectives for language understanding. Their LLM-JEPA framework showed that large language models benefit from predictive representation learning in ways that complement standard autoregressive training. \citet{liu2026jepareasoner} presented JEPA-Reasoner, decoupling latent-space reasoning from discrete token generation to enable unconstrained reasoning paths, while \citet{zhang2026thinkjepa} introduced ThinkJEPA, empowering latent world models with large vision-language reasoning capabilities. Collectively, these works indicate a growing trend toward unifying self-supervised prediction with language grounding, though none directly evaluate how these latent representations compare against instruction-tuned VLMs for discriminative video understanding tasks.


Beyond visual recognition, JEPA architectures serve as foundational world models for physical reasoning and decision-making. \citet{bagatella2025tdjepa} developed Temporal Difference JEPA, proving that latent-predictive representations enable zero-shot reinforcement learning across complex control environments. By combining temporal difference learning with JEPA's latent prediction, TD-JEPA demonstrated that forward predictive representations in a learned feature space encode sufficient dynamical information for effective policy transfer without environment-specific fine-tuning. \citet{terver2026physicaljepa} investigated the core drivers of success in physical planning using joint-embedding predictive world models, demonstrating strong generalization in physical simulation tasks. Their analysis identified that the latent prediction objective, rather than model scale alone, is the primary contributor to planning performance, providing evidence that JEPA representations capture causal-temporal structure that is qualitatively different from the associative representations learned through vision-language pretraining.


A fundamental distinction exists between the representation learning strategies of JEPA-based self-supervised models and instruction-tuned vision-language models. Self-supervised JEPA encoders such as V-JEPA~2.1 \citep{murlabadia2026vjepa21} learn visual representations through prediction in latent space, entirely without textual supervision. Their learned features are inherently visual-geometric, capturing motion trajectories, spatial configurations, and temporal dynamics as distributed patterns within a continuous embedding space. In contrast, vision-language models such as \qwenvl{} \citep{bai2025qwen2vl, qwen2025qwen3vl} acquire visual representations through cross-modal alignment with natural language, producing features that are linguistically grounded and semantically structured according to textual category boundaries. As discussed in the comprehensive survey by \citet{bordes2024vlm}, this linguistic grounding provides VLMs with strong zero-shot generalization capabilities, since visual concepts are inherently linked to compositional language descriptions. Recent efforts have sought to bridge the gap between linguistic and geometric representations; notably, \citet{lin2026qwen3d} introduced Qwen-3D, a generalist 3D vision-language model built upon the Qwen backbone that augments visual tokens with 3D Rotary Positional Embeddings and operates attention directly in 3D scene space, enabling scalable cross-view and temporal reasoning through explicit geometric awareness rather than purely linguistic alignment. However, \citet{zhang2023vlmsurvey} noted that vision-language alignment may bias representations toward semantically salient concepts that are easily described in natural language, potentially sacrificing fine-grained visual distinctions that lack clear linguistic labels.

The efficiency dimension further differentiates these paradigms. V-JEPA~2.1, with 654M parameters, achieves its representational power through architectural specialization and self-supervised pretraining on video data alone, whereas \qwenvl{} requires 7B parameters and massive multimodal corpora to achieve general-purpose visual understanding. \citet{pathak2026efficientvlms} surveyed quantization and pruning strategies for deploying large VLMs under hardware constraints, highlighting that the parameter efficiency of specialized self-supervised encoders represents a significant practical advantage for resource-constrained deployment scenarios.


While self-supervised JEPA models (\vjepaa{}, \citep{assran2025vjepa2}; V-JEPA 2.1, \citep{murlabadia2026vjepa21}) and large vision-language models (\qwenvl{}, \citep{qwen2025qwen3vl}) have demonstrated impressive individual capabilities, previous literature lacks a direct, controlled empirical comparison of their frozen-feature representations across fine-grained and coarse-grained video action recognition tasks. Furthermore, existing probing protocols rely on simple global average pooling, ignoring the rich multi-layer intermediate depth representations available in JEPA backbones. This study directly addresses these gaps by evaluating both representation paradigms across 8 baseline classifiers and introducing HAF-STD Net architecture to exploit multi-layer temporal dynamics.
